% !TeX root = théorie.tex

\section{Cardinaux}

\begin{proposition}
Soit $G = (S, A, p)$ un graphe de tas de sable.
$$
\left| \mathcal{M}\left(G\right) \right| = \prod_{s \in \widetilde{S}} d^+(s)
$$
\end{proposition}

\begin{proof}
Par définition de la stabilité, $\forall s \in \widetilde{S},
c(s) \in [\![0;d^+(s)[\![$.
\end{proof}

\begin{proposition}
Soit $G = (S, A, p)$ un graphe de tas de sable.
$$
\left| \mathcal{G}\left(G\right) \right| = \left| \det \widetilde{L} \right|
$$
\end{proposition}

\begin{proof}
Admis
\end{proof}

\begin{remarque}
Pour le graphe complet non orienté,
$\left|\mathcal{G}\left(K_n\right)\right| = n^{n+2}$.
Pour le graphe ligne non orienté,
$\left|\mathcal{G}\left(L_n\right)\right| = n+1$.
Pour les tas de sable de Kadanoff,
$\left|\mathcal{G}\left(\text{Ka}_{n, p}\right)\right| = p^n$.
(Calculs à expliciter)
\end{remarque}

\begin{remarque}
Pour tous les types de graphes que nous avons modélisé, il semble que le
quotient $\displaystyle \frac{\left|\mathcal{G}\right|}
{\left|\mathcal{M}\right|}$ tende vers 0
quand la taille du graphe tend vers l'infini.
\end{remarque}